% 作者来自sysu，如果想了解讨论数模比赛的更多内容和信息（例如ai提示词），可以访问个人主页https://github.com/sleepyy-dog/Survival-Handbook-by-KOKO

% 总体内容是和全国大学生数学建模高度相关的，在此之外还加入一点对于华数杯等中文赛的微调（国赛可以直接去掉）。

% 可能这个模板的话有点多，但是相信如果认真全都看下来对于数模新手和老手都会有帮助。

% 建议将此模板在overleaf或者texpage等在线编辑器编译

% 在vscode等本地ide中运行会遇到的一个问题是中文文件名不兼容，总之bug挺多，不建议

% 请使用 xelatex 来编辑！！！！！！！！！
\PassOptionsToPackage{quiet}{xeCJK}
\documentclass[withoutpreface,bwprint]{cumcmthesis}
\input{structure}

% 论文标题
 
\title{2026年全国大学生数学建模竞赛统一模板} 

%%%%%%%%%%%%%%%%%%%%%%%%%%%%%%%%%%%%%%%%%%%%%%%%%%%%%%%%%%%%%

\begin{document}

% 这个栏目模块是用于在一些要求的比赛例如华数杯，会要求你有此栏目，国赛不需要可以直接注释掉

\begin{table}
    \centering
    \renewcommand{\arraystretch}{1.8}  % 调整行高，1.8倍行距
    \begin{tabular}{|M{3cm}|M{8cm}|M{2.5cm}|}
        \hline
        \multicolumn{1}{|c|}{\large{所属类别}} & \multirow{2}{*}{\large\textbf{2026年"华数杯"数模竞赛(国赛注释此表)}} & \multicolumn{1}{|c|}{\large{参赛编号}}     \\ \cline{1-1} \cline{3-3} 
        \normalsize{本科组} & & \normalsize{123456} \\ 
        \hline
    \end{tabular}
\end{table}

\maketitle

\begin{abstract}

简单说明问题背景（一句浅一句深）。

\textbf{对于问题一,}（一两个句子总结整体方法、主要结果和结论）。

\textbf{对于问题二,}

\textbf{对于问题三,}

\textbf{对于问题四,}

（无总结）

\keywords{关键词;关键词;关键词;关键词;关键词(四到七个即可，五个差不多)}

\end{abstract}

%%%%%%%%%%%%%%%%%%%%%%%%%%%%%%%%%%%%%%%%%%%%%%%%%%%%%%%%%%%%%

\section{问题重述（这一部分总共就一页，配图之后可能有两页。）}

\subsection{问题背景（内容在7行左右）}

\subsection{问题提出}

前置介绍......

\textbf{问题1：}  

\textbf{问题2：}  

\textbf{问题3：} 

\textbf{问题4：}  

%%%%%%%%%%%%%%%%%%%%%%%%%%%%%%%%%%%%%%%%%%%%%%%%%%%%%%%%%%%%%

\section{问题分析（通常一页，加图片可能会有两页）}

\subsection{问题一分析}
对于问题一，

\subsection{问题二分析}	
对于问题二，

\subsection{问题三分析}
对于问题三，

\subsection{问题四分析}
对于问题四，

%%%%%%%%%%%%%%%%%%%%%%%%%%%%%%%%%%%%%%%%%%%%%%%%%%%%%%%%%%%%% 

\section{模型假设}

\textbf{（1）假设1}：

\textbf{（2）假设2}：

\textbf{（3）假设3}：

\textbf{（4）假设4}：

%%%%%%%%%%%%%%%%%%%%%%%%%%%%%%%%%%%%%%%%%%%%%%%%%%%%%%%%%%%%% 

\section{符号说明（参考条数为十到二十条）}
\begin{table}[H]
    \centering
    % \caption{符号说明} % 这个选择性增加
    \begin{tabularx}{\textwidth}{CLC}
        \toprule
        符号    & 说明    & 单位 \\
        \midrule
        $m     $& 质量 & $kg$ \\
        $V     $& 体积 & $m^3$ \\
        \bottomrule
    \end{tabularx}
    \label{tab:符号说明}
\end{table}

\section{问题一的模型的建立和求解}

有一个前置的流程介绍和流程图。

\subsection{模型建立}

\subsubsection{模型1名称}

详细流程，算法可画流程图和示意图。

\subsubsection{模型2名称}

具体内容
\subsection{模型求解结果}

不光呈现结果，要进行仿真绘图和数据绘图，和前文的假设或者直观感受相符。

还可以进行灵敏度的分析。
% %%%%%%%%%%%%%%%%%%%%%%%%%%%%%%%%%%%%%%%%%%%%%%%%%%%%%%%%%%%%% 

\section{问题二的模型的建立和求解}

\subsection{模型建立}

具体内容

\subsubsection{模型1名称}



\subsubsection{模型2名称}



\subsection{模型求解结果}

问题二结果及分析。

%%%%%%%%%%%%%%%%%%%%%%%%%%%%%%%%%%%%%%%%%%%%%%%%%%%%%%%%%%%%% 

\section{问题三的模型的建立和求解}

具体内容

\subsection{模型建立}



\subsubsection{模型1名称}



\subsubsection{模型2名称}



\subsection{模型求解结果}

问题三结果及分析。

%%%%%%%%%%%%%%%%%%%%%%%%%%%%%%%%%%%%%%%%%%%%%%%%%%%%%%%%%%%%% 

\section{问题四的模型的建立和求解}

\subsection{模型建立}

具体内容

\subsubsection{模型1名称}



\subsubsection{模型2名称}



\subsection{模型求解结果}

问题四结果及分析。

%%%%%%%%%%%%%%%%%%%%%%%%%%%%%%%%%%%%%%%%%%%%%%%%%%%%%%%%%%%%%

\section{模型的评价}

控制在三分之一页即可

\subsection{模型的优点（三个左右）}
\textbf{（1）优点1}：

\textbf{（2）优点2}：

\textbf{（3）优点3}：

\subsection{模型的缺点（两个或三个）}
\textbf{（1）缺点1}

\textbf{（2）缺点2}

%%%%%%%%%%%%%%%%%%%%%%%%%%%%%%%%%%%%%%%%%%%%%%%%%%%%%%%%%%%%%
%% 参考文献

% 这里对于参考文献的标注本来可以直接写在bib文件中来使用，但是用于这个文件无法自定义官方要求的ai使用的引用说明（DeepSeek,  DeepSeek-chat-v3.1,  深度求索(DeepSeek),  2025-09-06.），所以我们给出以下自定义解决的方案。

% {99} 表示序号最多为两位数，留出足够宽度
\begin{thebibliography}{99} 

% 遵循 GB/T 7714-2015
\bibitem{ref1} 昂温 G, 斯特朗 G. 概率论与数理统计[M]. 2版. 北京: 科学出版社, 2002.
\bibitem{ref2} 张三, 李四. 关于A算法的优化研究[J]. 计算机学报, 2022, 45(10): 2100-2110.
\bibitem{ai:deepseek1} DeepSeek,  DeepSeek-chat-v3.1,  深度求索(DeepSeek),  2025-09-06.
\end{thebibliography}

\newpage
%%%%%%%%%%%%%%%%%%%%%%%%%%%%%%%%%%%%%%%%%%%%%%%%%%%%%%%%%%%%%%% 附录
\begin{appendices}
\section{文件列表}
\begin{table}[H]
\centering
\begin{tabularx}{\textwidth}{CC}
\toprule
文件名   & 功能描述 \\
\midrule
P1.py & 问题一程序代码 \\
P2.py & 问题二程序代码 \\
P3.py & 问题三程序代码 \\
P4.py & 问题四程序代码 \\
\bottomrule
\end{tabularx}
\label{tab:文件列表}
\end{table}

\section{代码}
\noindent 

P1.py
\lstinputlisting[language=python]{code/q2.py}

P2.py
\lstinputlisting[language=python]{code/q2.py}

P3.py
\lstinputlisting[language=python]{code/q2.py}

P4.py
\lstinputlisting[language=python]{code/q2.py}
\end{appendices}
\end{document}